%% Le cube de 1 m de côté découpé en cubes de 1 dm : 10 x 10 x 10 = 1 000 dm3.
%% La conclusion chiffrée est la réponse : elle passe par \rappel.
\documentclass[tikz,border=4pt]{standalone}
\usepackage{liaisons}
\begin{document}
\begin{tikzpicture}[x=1cm,y=1cm,
  cote/.style={{Stealth[length=4pt,width=3.5pt]}-{Stealth[length=4pt,width=3.5pt]},
               line width=0.6pt, draw=black!65},
  etiq/.style={font=\scriptsize, text=black!75, inner sep=1.5pt}]

\def\a{3.6}      % côté du cube (cm sur le dessin)
\def\dx{1.35}    % fuite en x
\def\dy{0.95}    % fuite en y

%% faces : le cube vu en perspective cavalière
\fill[solideB!10] (0,0) -- (\a,0) -- (\a,\a) -- (0,\a) -- cycle;                       % face avant
\fill[solideB!18] (0,\a) -- (\dx,\a+\dy) -- (\a+\dx,\a+\dy) -- (\a,\a) -- cycle;       % dessus
\fill[solideB!6]  (\a,0) -- (\a+\dx,\dy) -- (\a+\dx,\a+\dy) -- (\a,\a) -- cycle;       % droite

%% quadrillage au pas de 1 dm (10 cases) sur les trois faces visibles
\foreach \i in {0,...,10} {
  \pgfmathsetmacro\t{\a*\i/10}
  \draw[line width=0.25pt, draw=solideB!45] (\t,0) -- (\t,\a);
  \draw[line width=0.25pt, draw=solideB!45] (0,\t) -- (\a,\t);
  \draw[line width=0.25pt, draw=solideB!35] (\t,\a) -- (\t+\dx,\a+\dy);
  \draw[line width=0.25pt, draw=solideB!35]
      ({\dx*\i/10},{\a+\dy*\i/10}) -- ({\a+\dx*\i/10},{\a+\dy*\i/10});
  \draw[line width=0.25pt, draw=solideB!35] (\a,\t) -- (\a+\dx,\t+\dy);
  \draw[line width=0.25pt, draw=solideB!35]
      ({\a+\dx*\i/10},{\dy*\i/10}) -- ({\a+\dx*\i/10},{\a+\dy*\i/10});
}

%% arêtes
\draw[line width=1pt, draw=solideB] (0,0) -- (\a,0) -- (\a,\a) -- (0,\a) -- cycle;
\draw[line width=1pt, draw=solideB] (0,\a) -- (\dx,\a+\dy) -- (\a+\dx,\a+\dy) -- (\a,\a);
\draw[line width=1pt, draw=solideB] (\a,0) -- (\a+\dx,\dy) -- (\a+\dx,\a+\dy);

%% cotes
\draw[cote] (0,-0.45) -- (\a,-0.45);
\node[etiq, anchor=north] at (\a/2,-0.5) {1 m = 10 dm};
\draw[cote] (-0.45,0) -- (-0.45,\a);
\node[etiq, anchor=east, rotate=90, anchor=south] at (-0.5,\a/2) {1 m = 10 dm};
\draw[cote] (\a+0.32,-0.22) -- (\a+\dx+0.32,\dy-0.22);
\node[etiq, anchor=north west] at (\a+\dx+0.3,\dy-0.3) {1 m};

%% la conclusion : c'est elle, la réponse
\node[anchor=north west, font=\small, inner sep=3pt] at (\a+\dx+0.9,\a+\dy)
  {$V = 1\,\mathrm{m}\times1\,\mathrm{m}\times1\,\mathrm{m}$};
\node[anchor=north west, font=\small, inner sep=3pt] at (\a+\dx+0.9,\a+\dy-0.5)
  {\rappel{$V = 10\,\mathrm{dm}\times10\,\mathrm{dm}\times10\,\mathrm{dm}$}};
\node[anchor=north west, font=\small\bfseries, text=solideA, inner sep=3pt]
  at (\a+\dx+0.9,\a+\dy-1.25) {\rappel{$1\ \mathrm{m^3} = 1\,000\ \mathrm{dm^3}$}};
\end{tikzpicture}
\end{document}
